\documentclass[11pt,a4paper]{article}

\usepackage[utf8]{inputenc}
\usepackage[T1]{fontenc}
\usepackage[spanish,es-tabla]{babel}
\usepackage{amsmath,amssymb,amsthm,mathtools}
\usepackage{booktabs}
\usepackage{enumitem}
\usepackage{geometry}
\usepackage{setspace}
\usepackage{hyperref}
\usepackage{microtype}

\geometry{margin=1in}
\onehalfspacing

\hypersetup{
  pdftitle={Nota formal: Conjunciones ordinales y sinergia irreducible},
  pdfauthor={Johel Padilla-Villanueva},
  pdfkeywords={RECD, Systemic Tau, conjuncion ordinal, sinergia irreducible, excess3}
}

% Theorems
\theoremstyle{definition}
\newtheorem{definition}{Definici\'on}[section]
\newtheorem{axiom}{Axioma}[section]
\newtheorem{remark}{Observaci\'on}[section]
\newtheorem{example}{Ejemplo}[section]
\theoremstyle{plain}
\newtheorem{proposition}{Proposici\'on}[section]
\newtheorem{lemma}{Lema}[section]
\newtheorem{conjecture}{Conjetura}[section]
\theoremstyle{remark}
\newtheorem{scholium}{Scholium}[section]

\newcommand{\Res}{\mathrm{Res}}
\newcommand{\RECD}{\mathrm{RECD}}
\newcommand{\TC}{\mathrm{TC}}
\newcommand{\Syn}{\mathrm{Syn}}
\newcommand{\Surp}{\mathrm{Surp}}
\newcommand{\excess}{\mathrm{excess3}}
\newcommand{\Rtwo}{\mathcal{R}_2}
\newcommand{\Jwin}{\mathcal{J}}
\newcommand{\claim}{\mathsf{Claim}}

\title{\textbf{Nota formal}\\[0.4em]
\large Tres niveles de conjunci\'on ordinal, anidamiento y\\
sinergia irreducible en el marco Systemic Tau / RECD\\[0.6em]
\normalsize (Puntos 1a--1c del roadmap te\'orico)}

\author{
Johel Padilla-Villanueva\\
\textit{Department of Environmental Health, University of Puerto Rico}\\
\texttt{johelpadilla@gmail.com}
}

\date{2026-08-01 \quad$\cdot$\quad v0.1}

\begin{document}
\maketitle

\begin{abstract}
Esta nota fija el \emph{objeto matem\'atico} de las conjunciones ordinales
en el paradigma Systemic Tau / RECD, separ\'andolo de los estimadores
operativos ($\Phi_1$, $\Phi_2$, $\excess$) y de la ponderaci\'on del reloj
($\alpha_\ell(\lambda)$). Se introducen tres axiomas m\'inimos de
conjunci\'on, se definen tres \emph{claims} estructurales de profundidad
creciente, se clarifican tres sentidos de ``anidamiento'' que no deben
confundirse, y se caracteriza la \textbf{sinergia irreducible} (Nivel~3)
como residuo del joint simb\'olico respecto de la clase de leyes de orden
$\le 2$. El proxy $\excess$ queda relegado a estimador de ese residuo, no
a su definici\'on. El resultado es un criterio n\'itido para responder:
\emph{qu\'e es} una sinergia irreducible en t\'erminos ordinales y
\emph{por qu\'e} el Nivel~3 no es ``m\'as correlaci\'on''.
\end{abstract}

\tableofcontents
\vspace{1em}

% ============================================================
\section{Prop\'osito, estatus y pureza}
\label{sec:purpose}

\subsection{Qu\'e cierra esta nota}

El paper de fundamentos \cite{Padilla2026framework} y la implementaci\'on
\texttt{nested-recd} definen estimadores operativos
$\Phi_1$, $\Phi_2$, $\excess$ y un acumulador
\[
\Delta\RECD(t)=\sum_{\ell=1}^{3}\alpha_\ell(\lambda)\,\Phi_\ell^\ast(t).
\]
All\'i el anidamiento $\Phi_3\supset\Phi_2\supset\Phi_1$ se declara
\emph{conceptual}, y la irreductibilidad del Nivel~3 se impone
estad\'isticamente v\'ia un proxy h\'ibrido, no por una definici\'on de
objeto independiente del estimador.

Esta nota suministra esa definici\'on. Sin ella, los puntos posteriores del
roadmap (PID, estatus de $\alpha(\lambda)$, Feigenbaum estructural,
invariancia) flotan sobre heur\'isticas.

\subsection{Tres planos que no se colapsan}
\label{sec:three_planes}

\begin{definition}[Tres planos del formalismo]
\label{def:planes}
\begin{enumerate}[label=(\roman*)]
  \item \textbf{Plano de claims} (objeto te\'orico): predicados o
  medidas abstractas $C_\ell$ / $M_\ell$ sobre la geometr\'ia relacional
  del proceso de s\'imbolos.
  \item \textbf{Plano de estimadores} (m\'etodo): $\Phi_1$, $\Phi_2$,
  $\excess$, $\Phi_3=\mathbf{1}\{\excess>\theta_3\}$.
  \item \textbf{Plano del reloj} (agregaci\'on): pesos
  $\alpha_\ell(\lambda)$ y el incremento $\Delta\RECD$.
\end{enumerate}
\end{definition}

\begin{scholium}[Pureza metaf\'isica]
Ninguna definici\'on de esta nota demuestra el \emph{acto de ser}, el
\emph{persistir} extramental, ni cierra conjeturas ontol\'ogicas de la
Summa (M-I-1). Las conjunciones son \emph{gram\'atica formal} de
organizaci\'on ordinal; el reloj RECD es un acumulador de esa gram\'atica,
no una prueba del ser.
\end{scholium}

\begin{scholium}[Pureza metodol\'ogica]
Declarar un estimador ``Nivel~3'' no implica autom\'aticamente
irreductibilidad. Solo el Plano~(i) fija qu\'e significa irreductible;
el Plano~(ii) estima; el Plano~(iii) pondera.
\end{scholium}

% ============================================================
\section{Universo de discurso}
\label{sec:universe}

\begin{definition}[Trayectoria multivariante]
Sea $d\ge 2$ y $X_t=\bigl(X_t^{(1)},\ldots,X_t^{(d)}\bigr)\in\mathbb{R}^d$
un proceso (o una realizaci\'on muestral finita). El caso $d=1$ queda
excluido: sin multivariante no hay conjunci\'on entre canales.
\end{definition}

\begin{definition}[Proceso de s\'imbolos Bandt--Pompe]
\label{def:symbols}
Fijados $m\ge 2$ y retardo $\tau\ge 1$ (defaults operativos: $m=3$,
$\tau=1$), sea $\pi_t^{(i)}\in S_m$ el patr\'on ordinal de la coordenada
$i$ en el vector de retardos
$\mathbf{v}_t^{(i)}=(X_t^{(i)},X_{t+\tau}^{(i)},\ldots,X_{t+(m-1)\tau}^{(i)})$,
con regla fija de desempate. El proceso de s\'imbolos es
\begin{equation}
\label{eq:S}
S_t=\bigl(\pi_t^{(1)},\ldots,\pi_t^{(d)}\bigr)\in\Sigma^d,
\qquad \Sigma:=S_m.
\end{equation}
Toda conjunci\'on de esta nota es funci\'on de $\{S_t\}$ (y de ventanas
finitas de $\{S_t\}$), no de la escala absoluta de $X$.
\end{definition}

\begin{definition}[Ley emp\'irica en ventana]
\label{def:J}
Para una longitud de ventana $w\ge 2$ y un tiempo $t\ge w-1$, sea
\begin{equation}
\Jwin_t^{(w)}
:=\frac{1}{w}\sum_{u=t-w+1}^{t}\delta_{S_u}
\end{equation}
la medida emp\'irica del joint en la ventana (con $\delta$ la masa de
Dirac en $\Sigma^d$). Cuando el contexto es claro escribimos $\Jwin$.
\end{definition}

% ============================================================
\section{Qu\'e es una conjunci\'on ordinal}
\label{sec:conjunction}

\subsection{Axiomas m\'inimos}

\begin{axiom}[Ordinalidad]
\label{ax:A1}
Una conjunci\'on es invariante bajo transformaciones componente a
componente estrictamente crecientes:
si $Z_t^{(i)}=f_i\bigl(X_t^{(i)}\bigr)$ con cada $f_i$ estrictamente
creciente, entonces los s\'imbolos de $Z$ coinciden con los de $X$ y
toda conjunci\'on definida solo sobre $S$ es id\'entica.
\end{axiom}

\begin{axiom}[Relacionalidad]
\label{ax:A2}
Una conjunci\'on depende de relaciones entre al menos dos coordenadas
(pares, hipergrafos o el joint), no de la din\'amica de un \'unico canal
en aislamiento.
\end{axiom}

\begin{axiom}[Profundidad]
\label{ax:A3}
El \'indice de nivel $\ell\in\{1,2,3\}$ clasifica la \emph{clase de
estructura exigida} (coincidencia; memoria relacional de pares;
residuo de orden superior), no meramente la magnitud de un score
escalar.
\end{axiom}

\begin{definition}[Conjunci\'on ordinal]
\label{def:conjunction}
Una \textbf{conjunci\'on ordinal} de nivel $\ell$ en el tiempo $t$ es un
par $(C_\ell(t), M_\ell(t))$ donde:
\begin{enumerate}[label=(\roman*)]
  \item $C_\ell(t)$ es un predicado (claim binario) sobre la geometr\'ia
  relacional de $S$ en $t$ (o en una ventana anclada en $t$);
  \item $M_\ell(t)\ge 0$ es una medida de evidencia del claim
  (score continuo o tasa normalizada);
  \item el par satisface los Axiomas~\ref{ax:A1}--\ref{ax:A3}.
\end{enumerate}
Los estimadores $\Phi_\ell$ del paper de fundamentos son instancias del
Plano~(ii) de $M_\ell$ (y, en el caso binario del Nivel~3, de $C_3$).
\end{definition}

\begin{remark}
Esta definici\'on impide identificar ``conjunci\'on'' con ``cualquier
correlaci\'on alta''. La correlaci\'on de Pearson sobre amplitudes puede
violar A1; un indicador univariante viola A2; un umbral sobre $|\tau_s|$
sin distinguir clase de estructura viola A3.
\end{remark}

% ============================================================
\section{Tres claims estructurales}
\label{sec:claims}

\subsection{Nivel 1 --- Coincidencia (sincron\'ia simb\'olica instant\'anea)}

\begin{definition}[Grafo de igualdad y Claim-1]
\label{def:L1}
El grafo de igualdad en $t$ es el grafo no dirigido $G_t^{\mathrm{eq}}$
sobre $\{1,\ldots,d\}$ con arista $ij$ sii $\pi_t^{(i)}=\pi_t^{(j)}$.
\begin{align}
\label{eq:C1}
\claim_1(t)
&:\iff E\bigl(G_t^{\mathrm{eq}}\bigr)\ne\emptyset
\quad\text{(o una fracci\'on m\'inima de aristas, seg\'un protocolo),}\\
\label{eq:M1}
M_1(t)
&:=\Phi_1(t)
=\frac{2}{d(d-1)}\sum_{1\le i<j\le d}
\mathbf{1}\{\pi_t^{(i)}=\pi_t^{(j)}\}\in[0,1].
\end{align}
\end{definition}

\begin{remark}[Tipo estructural]
Claim-1 afirma \emph{co-ocurrencia instant\'anea de s\'imbolos}. No
contiene memoria ni estructura de orden $>1$ m\'as all\'a de la igualdad
simult\'anea.
\end{remark}

\subsection{Nivel 2 --- Relaci\'on persistente (memoria en el grafo de pares)}

\begin{definition}[C\'odigo relacional]
\label{def:Rcode}
Para cada par $(i,j)$ y tiempo $t$,
\begin{equation}
R_t^{ij}
=
\begin{cases}
\mathrm{EQ} & \text{si }\pi_t^{(i)}=\pi_t^{(j)},\\
\mathrm{GT} & \text{si }\pi_t^{(i)}>\pi_t^{(j)},\\
\mathrm{LT} & \text{si }\pi_t^{(i)}<\pi_t^{(j)},
\end{cases}
\end{equation}
donde $>$ en $S_m$ denota un orden total fijo en el alfabeto de
permutaciones (p.ej.\ orden lexicogr\'afico del c\'odigo entero).
Alternativamente, el estado relacional puede tomarse como el par
ordenado $(\pi_t^{(i)},\pi_t^{(j)})$ sin colapsar a tres c\'odigos.
\end{definition}

\begin{definition}[Persistencia y Claim-2]
\label{def:L2}
Fijados $p_{\min}\ge 2$ y $f_{\min}\in(0,1]$ (defaults: $4$ y $0.75$),
el par $(i,j)$ es \emph{persistente} en $t$ si la misma relaci\'on
$R_t^{ij}$ ocupa al menos una fracci\'on $f_{\min}$ de la ventana
retrospectiva $\{t-p_{\min}+1,\ldots,t\}$. Entonces
\begin{align}
\claim_2(t)
&:\iff \text{existe al menos un par persistente en $t$},\\
M_2(t)
&:=\Phi_2(t)
=\frac{2}{d(d-1)}\sum_{i<j}
w(R_t^{ij})\,\mathbf{1}\{\text{$(i,j)$ persistente en $t$}\},
\end{align}
con $w\ge 0$ (p.ej.\ $w(\mathrm{EQ})=1$, $w(\mathrm{GT})=w(\mathrm{LT})=0.85$).
\end{definition}

\begin{remark}[Tipo estructural]
Claim-2 afirma \emph{organizaci\'on durable en el grafo de relaciones de
pares}. A\~nade memoria y calidad relacional; sigue viviendo en
dependencias de orden $\le 2$.
\end{remark}

\subsection{Nivel 3 --- Residuo irreducible (orden superior)}
\label{sec:L3}

Aqu\'i se fija el objeto \emph{antes} del proxy.

\begin{definition}[Clase de reducciones de orden $\le 2$]
\label{def:R2}
Sea $\mathcal{P}(\Sigma^d)$ el simplex de probabilidades sobre
$\Sigma^d$. La clase $\Rtwo\subset\mathcal{P}(\Sigma^d)$ consta de todas
las leyes $Q$ que se factorizan por m\'argenes y/o dependencias de a lo
sumo pares. Ejemplos can\'onicos (no exclusivos):
\begin{enumerate}[label=(\alph*)]
  \item \textbf{Independencia total:}
  $Q(s)=\prod_{i=1}^d q_i(s_i)$.
  \item \textbf{Modelos de pares (pairwise MRF / log-lineales de orden
  $\le 2$):}
  $Q(s)\propto\exp\bigl(\sum_i\theta_i(s_i)+\sum_{i<j}\theta_{ij}(s_i,s_j)\bigr)$.
  \item \textbf{Copias de marginales bipolares:}
  cualquier $Q$ cuyas proyecciones a subconjuntos de tama\~no $\ge 3$
  est\'en determinadas por las bipolares $\{Q_{ij}\}_{i<j}$.
\end{enumerate}
La elecci\'on operativa de $\Rtwo$ debe declararse en cada protocolo;
la definici\'on de irreductibilidad es relativa a esa clase.
\end{definition}

\begin{definition}[Residuo de orden superior]
\label{def:Res}
Sea $D(\,\cdot\,\|\,\cdot\,)$ una divergencia no negativa que se anula
sii sus argumentos coinciden (p.ej.\ KL, o una divergencia $f$
estricta). El \textbf{residuo de orden superior} de $\Jwin$ respecto de
$\Rtwo$ es
\begin{equation}
\label{eq:Res}
\Res\bigl(\Jwin;\Rtwo\bigr)
:=\inf_{Q\in\Rtwo}\, D\bigl(\Jwin\,\|\,Q\bigr)\in[0,+\infty].
\end{equation}
Si el \'infimo no se alcanza, se admite el \'infimo; en alfabetos finitos
y $\Rtwo$ cerrada en topolog\'ia adecuada, el m\'inimo existe para KL bajo
condiciones est\'andar de modelos exponenciales de pares.
\end{definition}

\begin{definition}[Claim-3: sinergia irreducible ordinal]
\label{def:L3}
Fijado un umbral $\theta_3\ge 0$ (posiblemente $0$ en el l\'imite
te\'orico),
\begin{align}
\label{eq:C3}
\claim_3(t)
&:\iff \Res\bigl(\Jwin_t^{(w)};\Rtwo\bigr)>\theta_3,\\
\label{eq:M3}
M_3(t)
&:=\Res\bigl(\Jwin_t^{(w)};\Rtwo\bigr).
\end{align}
Se dice que hay \textbf{sinergia irreducible (Nivel~3)} en $t$ cuando
$\claim_3(t)$ es verdadero.
\end{definition}

\begin{remark}[Lectura estructural]
Nivel~3 no afirma ``hay mucha dependencia''. Afirma: \emph{existe
estructura en el joint simb\'olico que no puede ser reproducida por
ninguna ley de la clase de orden $\le 2$}.
\end{remark}

\begin{definition}[Forma por proyecci\'on / invariante que muere]
\label{def:invariant_dies}
Sea $\Pi_2$ cualquier proyecci\'on que retiene solo informaci\'on de
pares (p.ej.\ el mapa $\Jwin\mapsto\{Q_{ij}\}_{i<j}$ seguido de un
m\'aximo-entrop\'ia en $\Rtwo$). Un funcional $I$ es un
\textbf{invariante de orden superior} si
\begin{equation}
I(\Jwin)>0
\quad\text{y}\quad
I\bigl(\Pi_2(\Jwin)\bigr)=0
\end{equation}
para la clase de $\Jwin$ de inter\'es. Entonces $I(\Jwin)>0$ implica
$\Res(\Jwin;\Rtwo)>0$ siempre que $I$ se anule en $\Rtwo$.
\end{definition}

Esta formulaci\'on es la v\'ia m\'as estable te\'oricamente: la sinergia
irreducible es lo que \emph{desaparece al proyectar a pares}.

% ============================================================
\section{Anidamiento, composici\'on y dependencia}
\label{sec:nesting}

\subsection{Tres sentidos de ``anidamiento'' (no intercambiables)}

\begin{definition}[Tres tipos de anidamiento]
\label{def:nesting_types}
\begin{enumerate}[label=(\Roman*)]
  \item \textbf{Anidamiento l\'ogico de eventos:}
  $\claim_3(t)\Rightarrow\claim_2(t)\Rightarrow\claim_1(t)$.
  \item \textbf{Anidamiento informacional / de claims:}
  el \emph{contenido estructural} exigido por el nivel $\ell{+}1$ es
  estrictamente m\'as fuerte (no reducible al de nivel $\ell$).
  \item \textbf{Anidamiento ontol\'ogico de reloj:}
  el peso $\alpha_\ell(\lambda)$ asigna masa creciente a niveles m\'as
  profundos bajo r\'egimen $\lambda$ alto (Plano~iii).
\end{enumerate}
\end{definition}

\begin{proposition}[Separaci\'on de tipos]
\label{prop:separation}
Con las definiciones operativas est\'andar de $\Phi_1$, $\Phi_2$ y
$\Res$:
\begin{enumerate}[label=(\roman*)]
  \item El Tipo~I \textbf{no} se cumple en general:
  puede haber $\Res>0$ sin pares persistentes, y pares persistentes sin
  igualdad instant\'anea dominante.
  \item El Tipo~II \textbf{s\'i} es el anidamiento can\'onico del
  paradigma: Claim-3 $\succ$ Claim-2 $\succ$ Claim-1 como
  \emph{exigencia estructural}.
  \item El Tipo~III es independiente de I--II: es dise\~no (o, en el
  punto~3 del roadmap, posible derivaci\'on) de $\alpha_\ell(\lambda)$.
\end{enumerate}
\end{proposition}

\begin{proof}[Esbozo]
(i)~$\Phi_1$ y $\Phi_2$ son funcionales de proyecciones de pares (con o
sin memoria); $\Res$ puede ser positivo por dependencias de ternas aun
cuando ninguna arista de $G_t^{\mathrm{eq}}$ sea densa y ninguna
relaci\'on sea persistente bajo $(p_{\min},f_{\min})$. Rec\'iprocamente,
un grafo de pares perfectamente persistente puede ser un modelo de
pares exacto ($\Res=0$). (ii)~Por construcci\'on de
Defs.~\ref{def:L1}--\ref{def:L3}. (iii)~Los pesos no alteran los claims.
\end{proof}

\begin{scholium}[Correcci\'on al s\'imbolo $\Phi_3\supset\Phi_2\supset\Phi_1$]
La inclusi\'on del paper de fundamentos debe leerse como \textbf{Tipo~II}
(anidamiento de claims), no como Tipo~I (implicaci\'on de indicadores).
Operativamente los estimadores se calculan en paralelo y se agregan por
\eqref{eq:drecd_ref}; la irreductibilidad se impone por $\Res$, no por
implicaci\'on l\'ogica de $\Phi_\ell$.
\end{scholium}

\subsection{Composici\'on del reloj}

\begin{definition}[Agregaci\'on RECD]
\label{def:composition}
\begin{equation}
\label{eq:drecd_ref}
\Delta\RECD(t)
=\sum_{\ell=1}^{3}\alpha_\ell(\lambda)\, \widehat{M}_\ell(t),
\end{equation}
donde $\widehat{M}_\ell$ es el estimador declarado de $M_\ell$ (p.ej.\
$\widehat{M}_1=\Phi_1$, $\widehat{M}_2=\Phi_2$,
$\widehat{M}_3=\excess$ o $\Res$). Esto es \textbf{composici\'on
ponderada de evidencias de claims distintos}, no uni\'on de
subconjuntos de eventos.
\end{definition}

\subsection{Dependencias entre niveles}

\begin{proposition}[Mapa de dependencias]
\label{prop:deps}
\begin{enumerate}[label=(\roman*)]
  \item $\claim_1\not\Rightarrow\claim_2$ (coincidencia sin persistencia).
  \item $\claim_2\not\Rightarrow\claim_3$ (pares persistentes con
  $\Res=0$).
  \item $\claim_3\Rightarrow$ ``existe estructura joint fuera de
  $\Rtwo$'' (por definici\'on).
  \item $\claim_3\not\Rightarrow\claim_2$ ni $\claim_1$ bajo los
  estimadores est\'andar.
\end{enumerate}
Si un protocolo exige anidamiento Tipo~I, debe \emph{redise\~nar} el
evento de Nivel~3 como
\begin{equation}
\label{eq:forced_nest}
\claim_3^{\mathrm{nest}}(t)
:\iff
\Res\bigl(\Jwin_t^{(w)};\Rtwo\bigr)>\theta_3
\;\wedge\;
\claim_2(t)
\end{equation}
(y opcionalmente $\wedge\,\claim_1(t)$). Esa es una decisi\'on de
dise\~no, no un teorema.
\end{proposition}

% ============================================================
\section{Por qu\'e el Nivel~3 no es ``m\'as correlaci\'on''}
\label{sec:not_more_corr}

\begin{definition}[Familia de observables de orden $\le 2$]
\label{def:order2_obs}
Llamamos observables de orden $\le 2$ a cualquier funcional de
$\{S_t\}$ que dependa solo de:
\begin{enumerate}[label=(\alph*)]
  \item m\'argenes univariantes y/o
  \item leyes bipolares (incluyendo coincidencias, c\'odigos $R^{ij}$,
  $\tau_K$ entre rangos, MI de pares, y persistencia de pares).
\end{enumerate}
En particular, $\Phi_1$, $\Phi_2$ y el Systemic Tau
\[
\tau_s(t)
=\frac{2}{d(d-1)}\sum_{i<j}\tau_K\bigl(r^{(i)}_t,r^{(j)}_t\bigr)
\]
(en ventanas) son de orden $\le 2$.
\end{definition}

\begin{proposition}[Separaci\'on Nivel~3 / correlaci\'on]
\label{prop:not_corr}
\begin{enumerate}[label=(\roman*)]
  \item Todo observable de orden $\le 2$ es constante en las fibras de
  $\Pi_2$ (misma estructura de pares $\Rightarrow$ mismo valor).
  \item $\Res(\,\cdot\,;\Rtwo)$ no es, en general, un observable de
  orden $\le 2$: puede variar con el joint manteniendo fijas todas las
  bipolares.
  \item Por tanto, el Nivel~3 no es una reescala de ``m\'as
  correlaci\'on'': es un residuo \emph{ortogonal} (en el sentido de
  proyecci\'on a $\Rtwo$) a la familia de observables de orden $\le 2$.
\end{enumerate}
\end{proposition}

\begin{example}[Esquema XOR ordinal]
\label{ex:xor}
Para $d=3$ y alfabeto adecuado, existen leyes $\Jwin$ tales que cada par
es independiente (o casi) mientras el joint est\'a concentrado en un
v\'inculo de tipo XOR/paridad. Entonces $\Phi_1$ y las MI de pares son
bajas, pero $\Res>0$. Sim\'etricamente, tres canales fuertemente
acoplados por un driver com\'un pueden tener MI de pares altas y
$\Res\approx 0$ tras ajustar un modelo de pares.
\end{example}

\begin{scholium}
La frase operativa correcta no es ``el Nivel~3 sube cuando hay m\'as
sincron\'ia'', sino: ``el Nivel~3 sube cuando el joint simb\'olico se
aleja de toda ley de orden $\le 2$''. En reg\'imenes ca\'oticos de mapas
acoplados ambos fen\'omenos pueden co-ocurrir; eso es una predicci\'on
emp\'irica, no la definici\'on.
\end{scholium}

% ============================================================
\section{Estatus de excess3 respecto de Res}
\label{sec:excess3}

El paper de fundamentos usa
\begin{equation}
\label{eq:excess3}
\excess(t)=0.6\,\Syn(t)+0.4\,\Surp(t),
\end{equation}
con
\begin{align}
\TC(t)
&=\sum_{i=1}^d H\bigl(\pi^{(i)}\bigr)
-H\bigl(\pi^{(1)},\ldots,\pi^{(d)}\bigr),\\
\Syn(t)
&=\bigl[\TC(t)-I_{\mathrm{pair}}(t)\bigr]_+,\\
\Surp(t)
&=\sum_s p_{\mathrm{obs}}(s)\,
\bigl[\log_2 p_{\mathrm{obs}}(s)-\log_2 p_{\mathrm{ind}}(s)\bigr]_+,
\end{align}
estimados en ventana (pesos $0.6/0.4$ fijos \emph{a priori}).

\begin{definition}[Estatus formal de $\excess$]
\label{def:excess_status}
$\excess$ es un \textbf{estimador-proxy} de $\Res(\Jwin;\Rtwo)$ (o de
una cota inferior de una variante de $\Res$), \emph{no} la definici\'on
de sinergia irreducible. En particular:
\begin{enumerate}[label=(\roman*)]
  \item no es un \'atomo axiom\'atico de PID completo;
  \item la combinaci\'on lineal $0.6/0.4$ es un contrato metodol\'ogico
  de falsabilidad, no un teorema;
  \item $\Phi_3=\mathbf{1}\{\excess>\theta_3\}$ es un test del claim,
  no el claim.
\end{enumerate}
\end{definition}

\begin{conjecture}[Puente proxy--residuo; agenda del punto~2]
\label{conj:bridge}
Existen constantes $c,C>0$ y una subclase $\mathcal{J}_\ast$ de leyes
emp\'iricas (p.ej.\ $d$ peque\~no, $m=3$, ventanas $w$ suficientes) tales
que
\begin{equation}
c\,\Res\bigl(\Jwin;\Rtwo^{\mathrm{pair}}\bigr)
\le \excess(\Jwin)
\le C\bigl(\Res\bigl(\Jwin;\Rtwo^{\mathrm{ind}}\bigr)+o(1)\bigr)
\end{equation}
bajo hip\'otesis de concentraci\'on expl\'icitas, donde
$\Rtwo^{\mathrm{ind}}$ es independencia total y $\Rtwo^{\mathrm{pair}}$
es la clausura de modelos de pares. Demostrar, refinar o delimitar esta
conjetura es el n\'ucleo del punto~2 del roadmap.
\end{conjecture}

\begin{remark}[Rutas del punto~2]
\begin{enumerate}[label=(\alph*)]
  \item \textbf{PID:} si un \'atomo sin\'ergico $S$ est\'a bien definido
  sobre $\Sigma^d$, tomar $M_3:=S$ y tratar $\excess$ como proxy de $S$.
  \item \textbf{O-information / interacciones de orden $k$:}
  identificar $\Res$ con interacci\'on de orden $\ge 3$.
  \item \textbf{Invariantes que mueren al proyectar:}
  Def.~\ref{def:invariant_dies} como definici\'on primaria (recomendada).
\end{enumerate}
\end{remark}

% ============================================================
\section{Respuestas n\'itidas al criterio de \'exito (parcial)}
\label{sec:answers}

\begin{enumerate}[label=\textbf{Q\arabic*.}, leftmargin=*]
  \item \textbf{¿Qu\'e es exactamente una sinergia irreducible en
  t\'erminos ordinales?}\\
  Es la validez de $\claim_3$: el residuo
  $\Res(\Jwin_t^{(w)};\Rtwo)>\theta_3$ del joint de s\'imbolos
  Bandt--Pompe respecto de la clase de leyes de orden $\le 2$.
  Equivalente: existencia de un invariante de orden superior que se
  anula al proyectar a pares.

  \item \textbf{¿Por qu\'e el Nivel~3 no es simplemente ``m\'as
  correlaci\'on''?}\\
  Porque correlaci\'on, $\tau_s$, $\Phi_1$ y $\Phi_2$ son observables de
  orden $\le 2$ (Def.~\ref{def:order2_obs}), mientras el Nivel~3 es el
  residuo respecto de $\Rtwo$ (Prop.~\ref{prop:not_corr}). Pueden moverse
  en direcciones opuestas (Ej.~\ref{ex:xor}).

  \item \textbf{¿Qu\'e cambia estructuralmente en las conjunciones cerca
  de una transici\'on cr\'itica?}\\
  \emph{A\'un no cerrado} (punto~4). Hip\'otesis de trabajo: la masa de
  $\Res$ y la fracci\'on $f_3$ del $\Delta\RECD$ crecen al cruzar la
  acumulaci\'on de Feigenbaum en mapas acoplados, de modo no reducible a
  un rescaling de $|\tau_s|$. La validaci\'on sint\'etica existente es
  evidencia, no teorema.

  \item \textbf{¿Qu\'e estatus tiene la ponderaci\'on ontol\'ogica
  $\alpha(\lambda)$?}\\
  \emph{A\'un no cerrado} (punto~3). En el estado actual es un
  \textbf{principio de dise\~no de r\'egimen} (Plano~iii) que eleva la
  masa de claims profundos cuando $\lambda$ se\~nala caos/alta
  volatilidad; no es a\'un una derivaci\'on desde un principio variacional
  \'unico. Esta nota lo separa limpia\-mente de la definici\'on de
  conjunci\'on.
\end{enumerate}

% ============================================================
\section{Agenda formal residual (puntos 2--4, 5--6)}
\label{sec:agenda}

\begin{enumerate}[label=\textbf{A\arabic*.}]
  \item \textbf{(Punto 2)} Demostrar o delimitar Conj.~\ref{conj:bridge};
  comparar $\Res$ con \'atomos PID y O-information en $\Sigma^d$;
  condiciones necesarias/suficientes para $\Res>0$ en mapas acoplados.
  \item \textbf{(Punto 3)} Estatus de $\alpha(\lambda)$: ¿par\'ametro
  fenomenol\'ogico, m\'axima entrop\'ia bajo constraints de claims,
  estabilidad de conjunciones, o consistencia ordinal a lo largo de
  bifurcaciones?
  \item \textbf{(Punto 4)} Caracterizar $\Res$, $M_\ell$ y $f_\ell$ a lo
  largo de la cascada de doblamiento de periodo: ¿qu\'e es invariante y
  qu\'e es caracter\'istico cerca de la acumulaci\'on de Feigenbaum?
  \item \textbf{(Punto 5)} Invariancia/covarianza de $\tau_s$ y de
  $(M_\ell)$ bajo reordenamientos, embeddings, cambios de escala y
  tama\~no muestral; versiones can\'onicas normalizadas.
  \item \textbf{(Punto 6)} Formulaci\'on categ\'orica opcional: claims
  como funtores desde la categor\'ia de procesos de s\'imbolos a
  posets de profundidad; $\Pi_2$ como funtor de olvido.
\end{enumerate}

% ============================================================
\section{Canon m\'inimo (para no reabrir)}
\label{sec:canon}

\begin{enumerate}[label=\textbf{C\arabic*.}]
  \item Conjunci\'on $:=$ claim + medida sobre $S$, sujeta a A1--A3.
  \item Tres planos: claim / estimador / reloj.
  \item Nivel~3 $:=$ $\Res(\Jwin;\Rtwo)>\theta_3$, no $:=\,\excess$.
  \item $\excess$ $:=$ proxy metodol\'ogico de $\Res$.
  \item Anidamiento can\'onico $:=$ Tipo~II (claims), no Tipo~I
  (indicadores), salvo redise\~no expl\'icito \eqref{eq:forced_nest}.
  \item $\Delta\RECD$ $:=$ suma ponderada de evidencias, no uni\'on
  conjuntista de eventos.
  \item M-I-1: nada de lo anterior demuestra el acto de ser.
\end{enumerate}

\vspace{1.5em}
\noindent\rule{\textwidth}{0.4pt}\\[0.5em]
\begin{center}
\em Fin de la Nota formal v0.1 (puntos 1a--1c).\\
Pr\'oximo entregable natural: puente $\Res$--$\excess$ / PID (punto 2a).
\end{center}

\begin{thebibliography}{9}
\bibitem{Padilla2026framework}
J.~Padilla-Villanueva.
\newblock Systemic Tau and the RECD Framework: A Relational Theory of
Hierarchical Ordinal Conjunctions and Critical Transitions in Complex
Systems.
\newblock Manuscrito de fundamentos, 2026.
\bibitem{BandtPompe2002}
C.~Bandt and B.~Pompe.
\newblock Permutation entropy: A natural complexity measure for time
series.
\newblock \emph{Phys.\ Rev.\ Lett.} 88, 174102 (2002).
\bibitem{Williams2010}
P.~L.~Williams and R.~D.~Beer.
\newblock Nonnegative decomposition of multivariate information.
\newblock arXiv:1004.2515 (2010).
\bibitem{Lizier2014}
J.~T.~Lizier, N.~M.~Timme, M.~Wibral, and J.~T.~Lizier.
\newblock (Referencias PID / informaci\'on de orden superior en el
marco de trabajo del paper de fundamentos.)
\end{thebibliography}

\end{document}
